\begin{landscape}
\begin{table}[p]
\centering
\small
\setlength{\tabcolsep}{5pt}
\renewcommand{\arraystretch}{1.25}

\begin{tabularx}{\linewidth}{
L{3.0cm}
L{3.1cm}
L{2.8cm}
L{3.1cm}
L{3.1cm}
Y
}
\hline
\textbf{Approach} &
\textbf{Automation Level} &
\textbf{Regrip Support} &
\textbf{Adaptability} &
\textbf{Runtime Integration} &
\textbf{Main Limitation} \\
\hline

Manual / teach pendant &
Low; poses and sequence are operator-defined. &
Manual only. &
Low; reteaching is required for changed geometry or placement. &
Direct or platform-specific execution. &
High programming effort and no integrated planning. \\
\hline

Offline programming tools &
Partial; simulation-assisted sequence preparation. &
Mostly manual or template-based. &
Moderate; model and sequence updates are still required. &
Tool-dependent. &
Expert sequence definition and repeated adjustment remain necessary. \\
\hline

\acs{cad}/digital twin &
Medium to high; model-based planning is possible. &
Possible when modelled explicitly. &
Moderate; depends on calibrated and updated models. &
External interfaces. &
Model preparation and calibration dominate deployment effort. \\
\hline

Academic planners &
High for isolated subproblems. &
Often strong in isolated regrasp planners. &
Narrow; usually scenario-specific. &
Usually not industrial runtime. &
Often not connected to an end-to-end industrial workflow. \\
\hline

Proposed approach &
High within the defined simulation scope. &
Automatic regrip decision and insertion. &
High for evaluated geometry and placement changes. &
drag\&bot + AnyRobot. &
Simulation validation and rectangular-geometry scope. \\
\hline

\end{tabularx}

\caption{Critical comparison of existing robotic programming and planning approaches with the proposed thesis approach.}
\label{tab:critical-comparison-existing-approaches}
\end{table}
\end{landscape}

\FloatBarrier